% GENERATED FILE. Edit canonical lessons, then run npm run generate:books.
% canonical-source-hash: fnv1a64:96c1ba599b764ecc
% canonical-lessons: PA-C233-aram, PA-C233-dari, PA-C233-vichhauna, PA-C233-khes, PA-C233-pakkha

\chapter{An Afternoon Rest}
\label{ch:pa-a2-an-afternoon-rest}
\hlchaptermodality{233}

\begin{chapteropening}
\textbf{By the end of this chapter:} \emph{I can talk about resting in the afternoon.}

Complete the last lesson of chapter 233: I can talk about resting in the afternoon.
\end{chapteropening}

\section[\texorpdfstring{\pa{ਆਰਾਮ} (ārām)}{ਆਰਾਮ (ārām)}]{\pa{ਆਰਾਮ} (ārām) — rest}

\label{lesson:PA-C233-aram}



\begin{quote}
Before the new one: say the Punjabi for to come unstitched, then the Punjabi for to stuff.
\end{quote}

\subsection*{You'll want to know: \pa{ਆਰਾਮ}}
\textbf{\pa{ਆਰਾਮ}} — \emph{ārām} — \textquotedblleft{}rest\textquotedblright{}.

\begin{quote}
\pa{ਦੁਪਹਿਰ} \pa{ਨੂੰ} \pa{ਦਾਦਾ} \pa{ਜੀ} \pa{ਆਰਾਮ} \pa{ਕਰਦੇ} \pa{ਹਨ।}
\end{quote}

\begin{quote}
\emph{dupahir nū̃ dādā jī ārām karde han.}
\end{quote}

\begin{quote}
My grandfather rests in the afternoon.
\end{quote}

A masculine noun.

\begin{cousinweb}
\textbf{\pa{ਆਰਾਮ}} (\emph{ārām}) is borrowed: it is usually traced to Persian \emph{ārām}, \textquotedblleft{}rest, quiet\textquotedblright{}.
\end{cousinweb}

\begin{grammarlens}[title={What you've built}]
One more word for saying what you do.
\end{grammarlens}

\subsection*{Your turn}
\begin{itemize}
\raggedright
  \item \emph{Say it:} \emph{ārām}
  \item \emph{Say it:} \emph{ārām}, once more
  \item \emph{Say it:} \emph{dupahir nū̃ dādā jī ārām karde han.} — My grandfather rests in the afternoon.
  \item \emph{Recall:} say the Punjabi for to come unstitched, then the Punjabi for to stuff, then say \emph{ārām} again
\end{itemize}

\begin{tcolorbox}[breakable,colback=teal!4,colframe=teal!35!black,title={Before you move on}]
What does \pa{ਆਰਾਮ} mean? (\textquotedblleft{}Rest\textquotedblright{}.) Say it once more.
\end{tcolorbox}
\section[\texorpdfstring{\pa{ਦਰੀ} (darī)}{ਦਰੀ (darī)}]{\pa{ਦਰੀ} (darī) — a cotton rug}

\label{lesson:PA-C233-dari}



\begin{quote}
Before the new one: say the Punjabi for to stuff, then the Punjabi for rest.
\end{quote}

\subsection*{You'll want to know: \pa{ਦਰੀ}}
\textbf{\pa{ਦਰੀ}} — \emph{darī} — \textquotedblleft{}a cotton rug\textquotedblright{}.

\begin{quote}
\pa{ਫ਼ਰਸ਼} \pa{ਉੱਤੇ} \pa{ਇੱਕ} \pa{ਨੀਲੀ} \pa{ਦਰੀ} \pa{ਹੈ।}
\end{quote}

\begin{quote}
\emph{farsh utte ikk nīlī darī hai.}
\end{quote}

\begin{quote}
There is a blue rug on the floor.
\end{quote}

A feminine noun.

\begin{cousinweb}
\textbf{\pa{ਦਰੀ}} (\emph{darī}) is inherited from a reconstructed \emph{*darita}, \textquotedblleft{}split cane, mat\textquotedblright{}; CDIAL lists Punjabi \emph{darī}, \textquotedblleft{}carpet\textquotedblright{} (CDIAL 6194).
\end{cousinweb}

\begin{grammarlens}[title={What you've built}]
One more word for saying what you do.
\end{grammarlens}

\subsection*{Your turn}
\begin{itemize}
\raggedright
  \item \emph{Say it:} \emph{darī}
  \item \emph{Say it:} \emph{darī}, once more
  \item \emph{Say it:} \emph{farsh utte ikk nīlī darī hai.} — There is a blue rug on the floor.
  \item \emph{Recall:} say the Punjabi for to stuff, then the Punjabi for rest, then say \emph{darī} again
  \item \emph{Write it:} copy \textbf{\pa{ਦਰੀ}} once, checking the whole word
\end{itemize}

\begin{tcolorbox}[breakable,colback=teal!4,colframe=teal!35!black,title={Before you move on}]
What does \pa{ਦਰੀ} mean? (\textquotedblleft{}A cotton rug\textquotedblright{}.) Say it once more.
\end{tcolorbox}
\section[\texorpdfstring{\pa{ਵਿਛਾਉਣਾ} (vichhāuṇā)}{ਵਿਛਾਉਣਾ (vichhāuṇā)}]{\pa{ਵਿਛਾਉਣਾ} (vichhāuṇā) — to spread out (a rug or bedding)}

\label{lesson:PA-C233-vichhauna}



\begin{quote}
Before the new one: say the Punjabi for rest, then the Punjabi for a cotton rug.
\end{quote}

\subsection*{You'll want to know: \pa{ਵਿਛਾਉਣਾ}}
\textbf{\pa{ਵਿਛਾਉਣਾ}} — \emph{vichhāuṇā} — \textquotedblleft{}to spread out (a rug or bedding)\textquotedblright{}.

\begin{quote}
\pa{ਮੈਂ} \pa{ਫ਼ਰਸ਼} \pa{ਉੱਤੇ} \pa{ਦਰੀ} \pa{ਵਿਛਾਉਂਦੀ} \pa{ਹਾਂ।}
\end{quote}

\begin{quote}
\emph{maiṁ farsh utte darī vichhāundī hāṁ.}
\end{quote}

\begin{quote}
I spread the rug on the floor.
\end{quote}

Past (I …): \pa{ਮੈਂ} \pa{ਚਾਦਰ} \pa{ਵਿਛਾਈ}, \emph{maiṁ chādar vichhāī}.

\begin{cousinweb}
\textbf{\pa{ਵਿਛਾਉਣਾ}} (\emph{vichhāuṇā}) is inherited from a reconstructed \emph{*vicchādayati}, \textquotedblleft{}spread out\textquotedblright{}; CDIAL lists Punjabi \emph{vichāuṇā}, \textquotedblleft{}to spread out (bedding \&c.)\textquotedblright{} (CDIAL 11692).
\end{cousinweb}

\begin{grammarlens}[title={What you've built}]
One more word for saying what you do.
\end{grammarlens}

\subsection*{Your turn}
\begin{itemize}
\raggedright
  \item \emph{Say it:} \emph{vichhāuṇā}
  \item \emph{Say it:} \emph{vichhāuṇā}, once more
  \item \emph{Say it:} \emph{maiṁ farsh utte darī vichhāundī hāṁ.} — I spread the rug on the floor.
  \item \emph{Recall:} say the Punjabi for rest, then the Punjabi for a cotton rug, then say \emph{vichhāuṇā} again
  \item \emph{Write it:} copy \textbf{\pa{ਵਿਛਾਉਣਾ}} once, checking the whole word
\end{itemize}

\begin{tcolorbox}[breakable,colback=teal!4,colframe=teal!35!black,title={Before you move on}]
What does \pa{ਵਿਛਾਉਣਾ} mean? (\textquotedblleft{}To spread out (a rug or bedding)\textquotedblright{}.) Say it once more.
\end{tcolorbox}
\section[\texorpdfstring{\pa{ਖੇਸ} (khes)}{ਖੇਸ (khes)}]{\pa{ਖੇਸ} (khes) — a thick cotton sheet}

\label{lesson:PA-C233-khes}



\begin{quote}
Before the new one: say the Punjabi for a cotton rug, then the Punjabi for to spread out.
\end{quote}

\subsection*{You'll want to know: \pa{ਖੇਸ}}
\textbf{\pa{ਖੇਸ}} — \emph{khes} — \textquotedblleft{}a thick cotton sheet\textquotedblright{}.

\begin{quote}
\pa{ਮੈਂ} \pa{ਮੰਜੇ} \pa{ਉੱਤੇ} \pa{ਖੇਸ} \pa{ਵਿਛਾਉਂਦਾ} \pa{ਹਾਂ।}
\end{quote}

\begin{quote}
\emph{maiṁ manje utte khes vichhāundā hāṁ.}
\end{quote}

\begin{quote}
I spread a thick cotton sheet on the cot.
\end{quote}

A masculine noun.

\begin{cousinweb}
\textbf{\pa{ਖੇਸ}} (\emph{khes}) is inherited from a reconstructed \emph{*kheśśa}, \textquotedblleft{}a particular sort of cloth\textquotedblright{}; CDIAL lists Punjabi \emph{khes}, \textquotedblleft{}a kind of cloth\textquotedblright{} (CDIAL 3924).
\end{cousinweb}

\begin{grammarlens}[title={What you've built}]
One more word for saying what you do.
\end{grammarlens}

\subsection*{Your turn}
\begin{itemize}
\raggedright
  \item \emph{Say it:} \emph{khes}
  \item \emph{Say it:} \emph{khes}, once more
  \item \emph{Say it:} \emph{maiṁ manje utte khes vichhāundā hāṁ.} — I spread a thick cotton sheet on the cot.
  \item \emph{Recall:} say the Punjabi for a cotton rug, then the Punjabi for to spread out, then say \emph{khes} again
  \item \emph{Write it:} copy \textbf{\pa{ਖੇਸ}} once, checking the whole word
\end{itemize}

\begin{tcolorbox}[breakable,colback=teal!4,colframe=teal!35!black,title={Before you move on}]
What does \pa{ਖੇਸ} mean? (\textquotedblleft{}A thick cotton sheet\textquotedblright{}.) Say it once more.
\end{tcolorbox}
\section[\texorpdfstring{\pa{ਪੱਖਾ} (pakkhā)}{ਪੱਖਾ (pakkhā)}]{\pa{ਪੱਖਾ} (pakkhā) — a fan}

\label{lesson:PA-C233-pakkha}



\begin{quote}
Before the new one: say the Punjabi for to spread out, then the Punjabi for a thick cotton sheet.
\end{quote}

\subsection*{You'll want to know: \pa{ਪੱਖਾ}}
\textbf{\pa{ਪੱਖਾ}} — \emph{pakkhā} — \textquotedblleft{}a fan\textquotedblright{}.

\begin{quote}
\pa{ਗਰਮੀ} \pa{ਵਿੱਚ} \pa{ਅਸੀਂ} \pa{ਪੱਖਾ} \pa{ਚਲਾਉਂਦੇ} \pa{ਹਾਂ।}
\end{quote}

\begin{quote}
\emph{garmī vich asīṁ pakkhā chalāunde hāṁ.}
\end{quote}

\begin{quote}
In the hot season we run the fan.
\end{quote}

A masculine noun.

\begin{cousinweb}
\textbf{\pa{ਪੱਖਾ}} (\emph{pakkhā}) is inherited from Sanskrit \emph{pakṣa}, \textquotedblleft{}wing, feather\textquotedblright{}; CDIAL lists Punjabi \emph{pakkhā}, \textquotedblleft{}fan\textquotedblright{} (CDIAL 7627).
\end{cousinweb}

\begin{grammarlens}[title={What you've built}]
One more word for saying what you do.
\end{grammarlens}

\subsection*{Your turn}
\begin{itemize}
\raggedright
  \item \emph{Say it:} \emph{pakkhā}
  \item \emph{Say it:} \emph{pakkhā}, once more
  \item \emph{Say it:} \emph{garmī vich asīṁ pakkhā chalāunde hāṁ.} — In the hot season we run the fan.
  \item \emph{Recall:} say the Punjabi for to spread out, then the Punjabi for a thick cotton sheet, then say \emph{pakkhā} again
  \item \emph{Write it:} copy \textbf{\pa{ਪੱਖਾ}} once, checking the whole word
\end{itemize}

\begin{tcolorbox}[breakable,colback=teal!4,colframe=teal!35!black,title={Before you move on}]
What does \pa{ਪੱਖਾ} mean? (\textquotedblleft{}A fan\textquotedblright{}.) Say it once more.
\end{tcolorbox}
